% Imporditud: tools/import_eksperimendid.py
% Allikas: Eesti füüsikaolümpiaadi eksperimendi ülesannete
%   kogu (Rael Kalda, 2023), ülesanne Ü64, lahendus L64.
\setAuthor{Erkki Tempel}
\setRound{lõppvoor}
\setYear{2014}
\setNumber{P E1}
\setDifficulty{2}
\setTopic{vedelike mehaanika}
\setVahendid{Plastikpulk, joonlaud, plastiliin, mahlajook (punane), puhas vesi, 20 \% suhkrulahus (roheline)}
\setPunktid{12}
\setAllikas{Eksperimendi kogu Ü64, lahendus lk 61}
\setLahendusseis{toores}

\prob{Tundmatu vedelik}
Määrake mahlajoogi suhkrusisaldus eeldusel, et joogi tihedus suureneb võrdeliselt suhkrusisalduse kasvuga.

\hint

\solu
Plastik torust ja plastiliinist saame teha areomeetri, pannes väikese tüki plastiliini ühte toru otsa. Areomeeter tuleb valmistada nii, et ta ujuks nii vees kui suhkrulahuses. Parima tulemuse saab siis, kui valmistatud areomeeter vajub vees võimalikult sügavale - veest jääb välja kuskil 2 mm pikkune pulk. Paigutades areomeetri vedelikku, saame mõõta selle, kui suur osa areomeetrist jääb vedelikust välja. Vette paigutades on veest välja ulatuva areomeetri toru pikkus $l_0$ ning 20 \% suhtkurlahuse korral $l_1$. Kuna lahuse tihedus (protsent) on lineaarses sõltuvuses veest välja ulatuva toru pikkusega (vastavalt jõudude tasakaalule), saame leida, et ühele pikkusühikule (lx) vastav protsent on

20 \% pl = $l_1$ $-$ $l_0$

Mõõtes ära selle, kui palju jääb toru välja (lmahl) mahlajoogi korral, saame leida ka mahlajoogi suhkrusisalduse p

p = (lmahl $-$ $l_0$) $\cdot$ pl

Korrektse lahenduse jaoks on vaja teostada vähemalt kolm kordusmõõtmist iga lahuse korral.
\probend
